% Generated draft: requires mathematical and source review.
\documentclass[11pt]{article}
\usepackage[margin=25mm]{geometry}
\usepackage{fontspec}
\setmainfont{Latin Modern Roman}
\usepackage{amsmath,amssymb,mathtools,unicode-math}
\setmathfont{Latin Modern Math}
\usepackage{enumitem,needspace,xurl,hyperref}
\setlength{\parindent}{0pt}
\setlength{\parskip}{5pt}
\setlength{\emergencystretch}{3em}
\newcommand{\N}{\mathbb{N}}
\newcommand{\Z}{\mathbb{Z}}
\newcommand{\Q}{\mathbb{Q}}
\newcommand{\R}{\mathbb{R}}
\newcommand{\C}{\mathbb{C}}
\newcommand{\F}{\mathbb{F}}
\newcommand{\A}{\mathbb{A}}
\newcommand{\PP}{\mathbb{P}}
\newcommand{\E}{\mathbb{E}}
\newcommand{\eps}{\varepsilon}
\newcommand{\dd}{\,\mathrm{d}}
\begin{document}
\section*{UnsolvedMath ID: 10400042}
\textbf{Source problem code:} catalogue entry
\subsection{Definitions and mathematical statement}
Let $R$ be a commutative ring with identity, in particular $\mathbb Z$ or $\mathbb Q$. For a set $\mathcal S$, write $R\mathcal S$ for its free $R$-module. Let $\mathcal K$ and $\mathcal M_K$ be the knot and associated-object sets used in the source, and retain its decreasing filtrations labeled loop and $\Delta\Delta$.

For every relevant index $l$, describe the following $R$-modules:
\[
\frac{F_l(R\mathcal M_K;\mathrm{loop})}{F_{l+1}(R\mathcal M_K;\mathrm{loop})},
\qquad
\frac{F_l(R\mathcal K;\Delta\Delta)}{F_{l+1}(R\mathcal K;\Delta\Delta)}.
\]
Also describe the image of the source's filtered map
\[
\Phi_l:F_l(R\mathcal K;\Delta\Delta)
\longrightarrow F_l(R\mathcal M_K;\mathrm{loop}).
\]
The third request concerns the image at the filtration level as printed, not merely an image on the two associated graded quotients.

The source extract does not define $M_K$, the ambient knot category, either local move, the filtration generators/index convention, or the ``above map.'' A map induced by an assignment $K\mapsto M_K$ could be extended $R$-linearly, but its actual construction and filtration compatibility must be recovered from the missing context. The spaces are modules for general $R$, not automatically vector spaces; rational ranks alone would not describe their integral structure. No surgery interpretation, move diagram, or generic finite-type filtration is substituted for the absent definitions.
\subsection{Short English statement}
Describe both associated graded move-filtration modules and the image of the source's filtered knot-to-object map, over a general unital commutative ring. The move, object, filtration, and map definitions are missing.
\subsection{Sources}
\begin{itemize}
\item Ohtsuki, \emph{Problems on invariants of knots and 3-manifolds} (2002), Problem 2.20, PDF page 44, \url{https://msp.org/gtm/2002/04/gtm-2002-04-024s.pdf}; accessed 2026-07-29, PDF-font-structure extraction.
\item Local registry \texttt{data/unsolvedmath/problems.json}, permanent UnsolvedMath ID $10400042$, source code catalogue entry, \emph{Problem 2.20}. Dataset attribution: Ulam AI and the UnsolvedMath Contributors.
\end{itemize}
\end{document}
